\section{The Conversion of Rene Guenon}

\label{sec:ConversionGuenon}

Did \textbf{Rene Guenon} ``convert''? In the article ``Conversions'', included in the collection \emph{Initiation and Spiritual Realization}, he explains that the word can be taken in two totally different senses. The first is an ``intellectual metamorphosis'', an ``interior transformation'', metanoia, or ``change of \emph{nous}'' that indicates the ``conscious passage of the ordinary and individual mind, normally turned toward sensible things, to its superior transposition.''

This is an interior transformation and has nothing to do with the second type which is commonly called a religious conversion. The second type is the exterior passage from one exoteric religious form to another; this type is just an exterior and contingent change. It often lacks ``any real importance'' and has ``nothing to do with pure spirituality.'' Guenon has ``little regard for such converts in general'' who are often associated with the ``most exaggerated sectarianism''. Guenon claims the opposite:

\begin{quotex}
Contrary to what takes place in `conversion', nothing here implies the attribution of the superiority of one traditional form over another. It is merely a question of what one might call reasons of spiritual expediency, which is altogether different from simple individual preference.
\end{quotex}


Nevertheless, even the first type of conversion may be associated with an outward change of form. However, that choice is not arbitrary and there may be valid reasons for it. Guenon elucidates:

\begin{quotex}
those who, for reasons of an esoteric and initiatic order, adopt a traditional form different from that to which they would seem to be linked by their origin do this either because their native tradition provides them with no possibility of an esoteric order, or because their chosen tradition gives them a foundation that is more appropriate to their nature, and consequently more favorable to their spiritual work.
\end{quotex}


So we see that for the elite, it is the intellectual conversion that is important. The outer form is secondary. Therefore, the various attempts to recover lost forms for their own sake or to critique living forms based solely on their exteriority are misguided. The criteria are these:

\begin{enumerate}


\item Is the form rich and varied enough to incorporate true metaphysical teachings?

\item Is the exterior form able to capture the allegiance of the general population, given their various capacities to understand?

\item Does the exterior form provide a safe haven within which the elite can accomplish their tasks?
\end{enumerate}



\flrightit{Posted on 2010-02-25 by Cologero}

\begin{center}* * *\end{center}

\begin{footnotesize}
\begin{sffamily}
\texttt{GFJF on 2010-02-26 at 00:01 said:}

I think the issue some take with Christianity lies in the third of your criteria. The Christian religion seems to have been the fountainhead of a large number of subversive movements -- though even as I write this I can anticipate your response will be that Christianity was Judaized. Well, fair enough, but wasn't it also paganized in an earlier period? It seems to me that what is of most value in historical Christianity is what is least essentially Christian about it.

\hfill

\end{sffamily}
\end{footnotesize}

